\subsection{Questão 50}

\subsubsection{Enunciado}

A refração da luz no material 2 depende, entre outros fatores, do índice de refração \(n_2\) do material 2.

A Fig.~33-52b mostra o ângulo de refração \(\theta_2\) em função de \(n_2\).

A escala do eixo vertical é definida por
\[
\theta_{2a} = 20{,}0^\circ
\qquad \text{e} \qquad
\theta_{2b} = 40{,}0^\circ.
\]

\begin{enumerate}
    \item[(a)] Qual é o índice de refração do material 1?

    \item[(b)] Se o ângulo de incidência aumenta para \(60^\circ\) e \(n_2 = 2{,}4\), qual é o valor de \(\theta_2\)?
\end{enumerate}

\begin{figure}[h]
    \centering
    \imagemquestao{figures/fig33-52.jpeg}
    \label{fig:33-52}
\end{figure}


\subsubsection{Resolução}

\textbf{(a) Determinação de \(n_1\)}

Pelo gráfico, para

\[
n_2=1{,}3,
\]

temos

\[
\theta_2=\theta_{2b}=40^\circ.
\]

Da figura, o ângulo de incidência é

\[
\theta_1 \approx 30^\circ.
\]

Aplicando a lei de Snell:

\[
n_1\sin\theta_1=n_2\sin\theta_2.
\]

Logo,

\[
n_1=\frac{n_2\sin\theta_2}{\sin\theta_1}
=\frac{(1{,}3)\sin40^\circ}{\sin30^\circ}.
\]

Assim,

\[
n_1\approx\frac{1{,}3(0{,}643)}{0{,}500}
\approx 1{,}67.
\]

Portanto,

\[
\boxed{n_1 \approx 1{,}67}.
\]

\bigskip

\textbf{(b) Novo ângulo de refração}

Para

\[
\theta_1=60^\circ
\qquad\text{e}\qquad
n_2=2{,}4,
\]

a lei de Snell fornece

\[
n_1\sin\theta_1=n_2\sin\theta_2.
\]

Substituindo os valores:

\[
(1{,}67)\sin60^\circ
=
(2{,}4)\sin\theta_2.
\]

\[
\sin\theta_2
=
\frac{1{,}67\sin60^\circ}{2{,}4}
\approx 0{,}602.
\]

Portanto,

\[
\theta_2=\sin^{-1}(0{,}602)
\approx 37^\circ.
\]

Logo,

\[
\boxed{\theta_2 \approx 37^\circ}.
\]

\vspace{1cm}
